\section{Flujo completo de la CNN: de la extracción de características a la clasificación}

\subsection{Panorama general de la arquitectura completa}

Una red CNN de clasificación completa puede dividirse en dos etapas:

\begin{figure}[H]
\centering
\includegraphics[width=0.9\textwidth]{slides-images/slide-51.jpg}
\caption{Flujo completo de la CNN: la primera mitad es el \textbf{aprendizaje de características} (repetición de convolución + ReLU + agrupamiento), la segunda mitad es la \textbf{clasificación} (aplanado + capa totalmente conectada + Softmax).\protect\footnotemark}
\end{figure}
\footnotetext{Diapositiva 51.}

\begin{enumerate}
    \item \textbf{Etapa de aprendizaje de características}:
    \begin{itemize}
        \item Aprender las características de la imagen de entrada mediante \textbf{convolución}
        \item Introducir no linealidad mediante \textbf{activación no lineal} (porque los datos del mundo real son no lineales)
        \item Reducir la dimensionalidad y preservar la invarianza espacial mediante \textbf{agrupamiento}
        \item Repetir el proceso anterior varias veces
    \end{itemize}
    \item \textbf{Etapa de clasificación}:
    \begin{itemize}
        \item \textbf{Aplanar} (flatten) el mapa de características final en un vector unidimensional
        \item Clasificar mediante una \textbf{capa totalmente conectada}
        \item Usar \textbf{Softmax} para obtener la probabilidad de cada clase
    \end{itemize}
\end{enumerate}

\subsection{Cabeza de clasificación: salida Softmax}

Después de que las capas de convolución y agrupamiento de la CNN producen una representación de características de alto nivel, esta se proyecta a las distintas clases mediante una capa totalmente conectada, y finalmente la función Softmax produce una distribución de probabilidad:

\begin{figure}[H]
\centering
\includegraphics[width=0.9\textwidth]{slides-images/slide-52.jpg}
\caption{Cabeza de clasificación de la CNN: las capas CONV y POOL producen características de alto nivel que, tras el aplanado, pasan por una capa totalmente conectada, y finalmente Softmax produce la probabilidad de cada clase.\protect\footnotemark}
\end{figure}
\footnotetext{Diapositiva 52.}

$$\text{softmax}(y_i) = \frac{e^{y_i}}{\sum_j e^{y_j}}$$

donde $y_i$ es la puntuación cruda (logit) de la $i$-ésima clase producida por la capa totalmente conectada; Softmax la convierte en un valor de probabilidad, y la suma de las probabilidades de todas las clases es igual a 1.

\subsection{Implementación en código: TensorFlow y PyTorch}

\begin{figure}[H]
\centering
\includegraphics[width=0.9\textwidth]{slides-images/slide-53.jpg}
\caption{Implementación de una CNN en TensorFlow/Keras: dos capas de convolución (32 y 64 filtros) + aplanado + capa totalmente conectada + salida Softmax de 10 clases.\protect\footnotemark}
\end{figure}
\footnotetext{Diapositiva 53.}

\begin{lstlisting}[caption={Implementación de CNN en TensorFlow/Keras}]
import tensorflow as tf

def generate_model():
    model = tf.keras.Sequential([
        # First convolutional layer
        tf.keras.layers.Conv2D(32, filter_size=3, activation='relu'),
        tf.keras.layers.MaxPool2D(pool_size=2, strides=2),

        # Second convolutional layer
        tf.keras.layers.Conv2D(64, filter_size=3, activation='relu'),
        tf.keras.layers.MaxPool2D(pool_size=2, strides=2),

        # Fully connected classifier
        tf.keras.layers.Flatten(),
        tf.keras.layers.Dense(1024, activation='relu'),
        tf.keras.layers.Dense(10, activation='softmax')  # 10 outputs
    ])
    return model
\end{lstlisting}

\begin{lstlisting}[caption={Implementación de CNN en PyTorch}]
import torch
import torch.nn as nn

def generate_model():
    model = nn.Sequential(
        # First and second convolutional layer
        nn.Conv2d(in_channels=3, out_channels=32, kernel_size=3),
        nn.ReLU(),
        nn.MaxPool2d(kernel_size=2, stride=2),

        nn.Conv2d(in_channels=32, out_channels=64, kernel_size=3),
        nn.ReLU(),
        nn.MaxPool2d(kernel_size=2, stride=2),

        # Fully connected classifier
        nn.Flatten(),
        nn.Linear(64*6*6, 1024),  # flattened dim after 2 conv layers
        nn.ReLU(),
        nn.Linear(1024, 10),      # 10 outputs
    )
    return model
\end{lstlisting}

\begin{knowledgebox}{El "arte y la ciencia" de elegir los hiperparámetros}
Amini reconoce abiertamente que la elección de los hiperparámetros de la CNN (número y tamaño de los filtros, número de capas, etc.) es "más un arte que una ciencia" (more art than science), pero también existen algunas pautas empíricas:
\begin{itemize}
    \item El tamaño de los filtros suele ser pequeño ($3 \times 3$ o $5 \times 5$), porque las características locales no requieren un campo receptivo grande
    \item El número de filtros suele \textbf{aumentar capa por capa} (por ejemplo, $32 \rightarrow 64 \rightarrow 128$), porque las capas más altas necesitan representar más tipos de características compuestas
    \item Las dimensiones espaciales \textbf{disminuyen capa por capa} mediante el agrupamiento, en complemento con el aumento del número de filtros
\end{itemize}
\end{knowledgebox}

\subsection{Resumen de la sección}

Una red CNN de clasificación completa se divide en dos etapas: aprendizaje de características y clasificación. El aprendizaje de características extrae características abstractas capa por capa mediante el apilamiento de convolución + ReLU + agrupamiento; la etapa de clasificación proyecta las características a probabilidades de clase mediante el aplanado y una capa totalmente conectada. Tanto TensorFlow como PyTorch ofrecen API sencillas para construir una CNN.

\newpage
